\documentclass[11pt,compress,t,notes=noshow, xcolor=table]{beamer}

\input{../../style/preamble}
\input{../../latex-math/basic-math}
\input{../../latex-math/basic-ml}
\input{../../latex-math/optim-basics}

\newcommand{\Lambdamat}{\bm{\Lambda}}

\title{Optimization in Machine Learning}

\begin{document}

\titlemeta{
Constrained Optimization
}{
Norm-Constrained Optimization
}{
figure/norm_constrained_circle.pdf
}{
\item Why norm constraints appear in machine learning
\item Sphere constraints, orthogonality, and eigenvector problems
\item Matrix Lagrangian stationarity and variational eigenvalue theorems
\item Smallest and largest invariant subspaces, with PCA links
}

\begin{framev}[fs=small]{Why Norm Constraints Matter}
\textbf{Common forms}
$$
\min_{\xv} f(\xv)
\qquad
\text{s.t. }
\qquad
\|\xv\|_2 \le r
$$
$$
\max_{\xv} \xv^\top \A \xv
\qquad
\text{s.t. }
\qquad
\|\xv\|_2 = 1
$$
\splitV[0.48]{
\begin{itemize}
\item A norm bound limits scale and keeps the search in a compact region.
\item A unit-norm equality removes arbitrary scaling, so the variable represents a \textbf{direction}.
\item In ML, this shows up in PCA, SVD, and spectral methods.
\end{itemize}
}{
\image[1]{figure/norm_constrained_circle.pdf}
}
\end{framev}


\begin{framei}[fs=small]{Single-Vector Case: Rayleigh Quotient}
\item Assume $\A \in \R^{d \times d}$ is symmetric.
\item For one direction, the core problem is
$$
\min_{\|\xv\|_2 = 1} \xv^\top \A \xv
\qquad \text{or} \qquad
\max_{\|\xv\|_2 = 1} \xv^\top \A \xv.
$$
\item Because $\|\xv\|_2 = 1$, the objective equals the \textbf{Rayleigh quotient}
$$
R_{\A}(\xv) := \frac{\xv^\top \A \xv}{\|\xv\|_2^2}.
$$
\item The feasible set is the unit sphere, so the search is over directions only.
\item The extremizers are eigenvectors of $\A$:
$$
\min_{\|\xv\|_2 = 1} \xv^\top \A \xv = \lambda_{\min}(\A),
\qquad
\max_{\|\xv\|_2 = 1} \xv^\top \A \xv = \lambda_{\max}(\A).
$$
\item The multi-vector problem is the orthogonal extension of this fact.
\end{framei}


\begin{framei}[fs=small]{Matrix Form for $k$ Orthogonal Directions}
\item Stack the directions into a matrix
$$
\Xmat := [\xv_1,\dots,\xv_k] \in \R^{d \times k}.
$$
\item Then
$$
\sumik \xv_i^\top \A \xv_i = \trace(\Xmat^\top \A \Xmat),
\qquad
\|\xv_i\|_2 = 1 \text{ and } \xv_i^\top \xv_j = 0
\iff
\Xmat^\top \Xmat = \id_k.
$$
\item The minimization problem becomes
$$
\min_{\Xmat \in \R^{d \times k}} \trace(\Xmat^\top \A \Xmat)
\qquad
\text{s.t. } \qquad
\Xmat^\top \Xmat = \id_k.
$$
\item Interpretation: choose an orthonormal $k$-dimensional subspace with minimum total quadratic energy.
\item Important caveat: the feasible set $\{\Xmat \mid \Xmat^\top \Xmat = \id_k\}$ is \textbf{nonconvex}.
\end{framei}


\begin{framei}[fs=small]{Matrix Lagrangian Stationarity}
\item Use one symmetric matrix multiplier $\Lambdamat \in \R^{k \times k}$ for the whole orthogonality constraint:
$$
\mathcal{L}(\Xmat,\Lambdamat)
:=
\trace(\Xmat^\top \A \Xmat)
-
\trace\!\big(\Lambdamat(\Xmat^\top \Xmat - \id_k)\big).
$$
\item Stationarity with respect to $\Xmat$ gives
$$
\nabla_{\Xmat} \mathcal{L}
=
2 \A \Xmat - 2 \Xmat \Lambdamat
=
\zero
\qquad \Longrightarrow \qquad
\A \Xmat = \Xmat \Lambdamat.
$$
\item This simultaneously handles unit norms and orthogonality; we no longer need one scalar multiplier per column.
\item Teaching point: this is better viewed as \textbf{Lagrangian stationarity plus spectral structure}, not as a loose scalar \enquote{dual derivation}.
\end{framei}


\begin{framei}[fs=small]{From Stationarity to Invariant Subspaces}
\item The relation
$$
\A \Xmat = \Xmat \Lambdamat
$$
\item means that the column space $\operatorname{span}(\Xmat)$ is an \textbf{invariant subspace} of $\A$.
\item Because $\A$ is symmetric, we can choose an orthonormal basis inside that subspace that diagonalizes $\Lambdamat$:
$$
\A \Xmat
=
\Xmat \operatorname{diag}(\lambda_{i_1},\dots,\lambda_{i_k}).
$$
\item In that basis, the columns of $\Xmat$ are eigenvectors of $\A$.
\item Any rotated basis $\Xmat \mathbf{Q}$ with $\mathbf{Q}^\top \mathbf{Q} = \id_k$ spans the same subspace and has the same objective value.
\end{framei}


\begin{framei}[fs=small]{Variational Characterization}
\item The multi-vector problem is solved by the \textbf{Ky Fan / Courant-Fischer principle}.
\item If $\lambda_1 \le \lambda_2 \le \dots \le \lambda_d$ are the eigenvalues of $\A$, then
$$
\min_{\Xmat^\top \Xmat = \id_k} \trace(\Xmat^\top \A \Xmat)
=
\sum_{i=1}^k \lambda_i,
$$
\[
\max_{\Xmat^\top \Xmat = \id_k} \trace(\Xmat^\top \A \Xmat)
=
\sum_{i=d-k+1}^d \lambda_i.
\]
\item This is the theorem-backed extension of the Rayleigh quotient from one vector to $k$ orthonormal vectors.
\item The theorem chooses an optimal \textbf{subspace}; the particular orthonormal basis inside that subspace is secondary.
\end{framei}


\begin{framei}[fs=small]{Result for the Minimization Variant}
\item Let $\lambda_1 \le \lambda_2 \le \dots \le \lambda_d$ be the eigenvalues of $\A$, with orthonormal eigenvectors $\vv_1,\dots,\vv_d$.
\item Any minimizer satisfies
$$
\operatorname{span}(\Xmat^\ast)
=
\operatorname{span}(\vv_1,\dots,\vv_k)
\qquad
\text{when } \lambda_k < \lambda_{k+1}.
$$
\item More generally, if eigenvalues are repeated, any orthonormal basis of the corresponding optimal invariant subspace is optimal.
\item The optimal value is
\[
p^\ast = \sum_{i=1}^k \lambda_i.
\]
\item So the key object is the optimal \textbf{subspace}, not one fixed ordered list of basis vectors.
\end{framei}


\begin{framev}[fs=small]{Maximization Variant}
\splitV[0.5]{
The maximization problem keeps the same feasible set,
\[
\max_{\Xmat^\top \Xmat = \id_k} \trace(\Xmat^\top \A \Xmat),
\]
but chooses the dominant invariant subspace.
\begin{itemize}
\item The optimizer picks the \textbf{largest $k$ eigenvalues}.
\item Any orthonormal basis of the span of the corresponding top-$k$ eigenvectors is optimal.
\item For $k=1$:
\[
\max_{\|\xv\|_2 = 1} \xv^\top \A \xv = \lambda_{\max}(\A).
\]
\end{itemize}
}{
\imageC[0.75]{figure/norm_constrained_spectrum.pdf}
}
\end{framev}


\begin{framev}[fs=small]{ML Connection: Principal Components}
\textbf{PCA}
$$
\max_{\Xmat^\top \Xmat = \id_k} \trace(\Xmat^\top \Sigma \Xmat)
$$
finds the principal subspace of largest variance.\\
\spacer
\splitV[0.45]{
\begin{itemize}
\item PC1 is the dominant eigenvector of the covariance matrix $\Sigma$.
\item Further components add orthogonality and therefore recover the top-$k$ eigenspace.
\item PCA is exactly a norm-constrained eigenvalue problem in the maximization form.
\end{itemize}
}{
\image[1]{figure/norm_constrained_pca.pdf}
}
\end{framev}


\begin{framei}[fs=small]{Why This Special Nonconvex Case Still Works}
\item The feasible set $\{\Xmat \mid \Xmat^\top \Xmat = \id_k\}$ is nonconvex, so the usual convex-duality story does \emph{not} apply automatically.
\item Exact solvability comes from special algebra:
\begin{itemize}
\item symmetry gives an orthonormal eigenbasis
\item stationarity forces an invariant subspace
\item the variational theorem identifies which subspace is optimal
\end{itemize}
\item \textbf{Spectral clustering:} typically uses the smallest nontrivial eigenvectors of a graph Laplacian to embed nodes before clustering.
\item \textbf{SVD:} dominant singular vectors solve norm-constrained extremal problems for $\mathbf{M}^\top \mathbf{M}$ and $\mathbf{M}\mathbf{M}^\top$.
\item This is why eigenvalue problems show up repeatedly across machine learning: the geometry and the linear algebra match perfectly.
\end{framei}


\begin{framev}{Take-Away}
\begin{enumerate}
\item Norm constraints often turn \enquote{find the best vector} into \enquote{find the best direction}.
\item In matrix form, orthogonality is the single constraint $\Xmat^\top \Xmat = \id_k$.
\item Stationarity gives an invariant subspace equation $\A \Xmat = \Xmat \Lambdamat$.
\item Variational eigenvalue theorems then identify the optimal bottom-$k$ or top-$k$ subspace.
\end{enumerate}
\vfill
\oneliner{Minimization picks the smallest eigenspace; maximization picks the largest one.}
\end{framev}

\endlecture
\end{document}
